% ==============================================================================
% Supplementary Materials
% "What Changes When LLMs Learn to Reason?
%  Mechanistic Signatures of Cognitive Bias Processing"
% ICML 2026 MechInterp Workshop
% ==============================================================================

\documentclass[accepted]{article}

% ---------- Packages (matched to main paper) ----------
\usepackage[utf8]{inputenc}
\usepackage[T1]{fontenc}
\usepackage{amsmath,amssymb,amsfonts}
\usepackage{graphicx}
\usepackage{booktabs}
\usepackage{enumitem}
\usepackage{hyperref}
\usepackage{url}
\usepackage{xcolor}
\usepackage{natbib}
\usepackage{microtype}
\usepackage[margin=1in]{geometry}
\usepackage{longtable}
\usepackage{array}

% ---------- Custom commands (matched to main paper) ----------
\newcommand{\sone}{S1-like}
\newcommand{\stwo}{S2-like}
\newcommand{\pzero}{\texttt{P0}}
\newcommand{\data}[1]{\textbf{[#1]}}

% ---------- Title ----------
\title{Supplementary Materials:\\
  What Changes When LLMs Learn to Reason?\\
  Mechanistic Signatures of Cognitive Bias Processing}

\author{%
  \textbf{Bright Liu}\textsuperscript{1} \quad
  \textbf{[Co-authors]}\textsuperscript{1} \quad
  \textbf{[Advisor]}\textsuperscript{1,2} \\[6pt]
  \textsuperscript{1}Harvard Undergraduate AI Safety (HUSAI) \quad
  \textsuperscript{2}[Department] \\
}

\date{}

\begin{document}

\maketitle

\noindent This document provides additional details on the benchmark
design, full behavioral and probing results, cross-prediction analyses,
transfer matrices, representational geometry, and lure susceptibility
distributions that complement the main text.  Section labels correspond
to the appendix identifiers referenced in the main paper.

% ============================================================
% A. Full Benchmark Description
% ============================================================

\section*{A. Full Benchmark Description}
\label{sec:appendix_benchmark}

\subsection*{A.1 Overview}

The benchmark comprises \textbf{470 items} organized into \textbf{11
categories} spanning \textbf{4 heuristic families}.  Of these, 380
items (190 matched conflict/control pairs) across 8 categories form the
primary evaluation set reported in the main text.  The remaining 90
items (45 pairs) belong to 3 additional categories added for extended
analysis: a natural frequency framing subset, sunk cost variants, and
contamination baselines.

Every conflict item has a structurally matched no-conflict control that
preserves surface features (narrative length, domain, vocabulary) while
removing the heuristic--normative tension.  This matched-pair design
follows De~Neys' conflict detection paradigm \citep{deneys2012bias}.

\subsection*{A.2 Design Principles}

\begin{enumerate}[leftmargin=*,itemsep=4pt]
  \item \textbf{Novel isomorphs.}  All items use novel surface features
  (names, professions, quantities, scenarios) to resist memorization
  from pre-training corpora.  Classic versions (e.g., the original
  bat-and-ball problem, the Linda problem) are included only as
  contamination baselines and excluded from all primary analyses.

  \item \textbf{Matched pairs.}  Each conflict item is paired with a
  structurally identical control.  For base rate neglect, the control
  aligns the diagnostic description with the base rate so no conflict
  exists.  For conjunction fallacy, the control presents the single
  event as more representative than the conjunction.  For syllogisms,
  the control aligns conclusion believability with logical validity.

  \item \textbf{Difficulty calibration.}  Items are assigned difficulty
  ratings from 1 (trivial) to 5 (maximal heuristic pull) based on pilot
  testing.  Difficulty ratings are balanced across conflict/control pairs
  within each category.

  \item \textbf{Natural frequency framing.}  A 10-pair subset of base
  rate neglect items presents statistical information in Gigerenzer's
  natural frequency format (``10 out of 1000'' rather than ``1\%'') to
  test the ecological rationality critique
  \citep{gigerenzer1995improve}.
\end{enumerate}

\subsection*{A.3 Category Table}

Table~\ref{tab:benchmark_full} provides the full category breakdown
with example prompts and sourcing information.

\begin{table}[ht]
\centering
\caption{Full benchmark description.  All items use novel surface
  features; classic versions are contamination baselines only.
  ``Freq.''\ = natural frequency framing subset.}
\label{tab:benchmark_full}
\vspace{4pt}
\small
\begin{tabular}{p{2.6cm}ccp{3.8cm}p{2.8cm}}
\toprule
\textbf{Category} & \textbf{\# Pairs} & \textbf{Heuristic Family} & \textbf{Example Prompt (abbrev.)} & \textbf{Source} \\
\midrule
Base rate neglect & 35 & Representativeness &
  ``In a group of 995 engineers and 5 lawyers, a person is described as analytical\ldots Is this person more likely a lawyer or engineer?'' &
  Kahneman \& Tversky (1973); novel isomorphs \\
\addlinespace
\quad Freq.\ subset & (10) & Representativeness &
  ``10 out of 1000 people have condition X. A test has 95\% sensitivity\ldots'' &
  Gigerenzer \& Hoffrage (1995) \\
\addlinespace
Conjunction fallacy & 20 & Representativeness &
  ``Mx.\ Rivera is 31, outspoken, studied philosophy\ldots Which is more probable: (A)~bank teller, (B)~bank teller and activist?'' &
  Tversky \& Kahneman (1983); novel isomorphs \\
\addlinespace
Belief-bias syllogisms & 25 & Belief bias &
  ``All flowers need water. Roses need water. Therefore, roses are flowers.'' (valid/invalid $\times$ believable/unbelievable) &
  Evans et al.\ (1983); novel isomorphs \\
\addlinespace
Sunk cost fallacy & 15 & Loss aversion &
  ``You have invested \$9M in Project Aurora. A competitor has launched a superior product\ldots Continue investing or cut losses?'' &
  Arkes \& Blumer (1985); novel isomorphs \\
\addlinespace
CRT variants & 30 & Cognitive reflection &
  ``A mug and a lid cost \$1.10 in total. The mug costs \$1.00 more than the lid\ldots'' &
  Frederick (2005); structural isomorphs \\
\addlinespace
Arithmetic & 25 & Computation &
  ``Calculate $247 \times 18$ step by step.'' (with/without carrying traps) &
  Original \\
\addlinespace
Framing effects & 20 & Loss aversion &
  ``Program A saves 200 people. Program B: 1/3 chance all 600 saved, 2/3 chance none saved.'' (gain/loss framing) &
  Tversky \& Kahneman (1981); novel scenarios \\
\addlinespace
Anchoring & 20 & Anchoring \& adjustment &
  ``The tallest building in the world is 828m. Estimate the average height of buildings in downtown Portland.'' (extreme vs.\ neutral anchor) &
  Tversky \& Kahneman (1974); novel anchors \\
\addlinespace
Classic baselines & --- & Mixed &
  Original bat-and-ball, Linda, etc.\ & Contamination check only \\
\bottomrule
\end{tabular}
\end{table}

\subsection*{A.4 Scoring}

Each item is scored on a four-way classification:
\begin{itemize}[leftmargin=*,itemsep=2pt]
  \item \textbf{Correct}: the normatively correct response.
  \item \textbf{S1-lure}: the specific heuristic-predicted wrong answer
  (e.g., ``lawyer'' when base rates strongly favor ``engineer'').
  \item \textbf{Other-wrong}: an incorrect response that does not match
  the heuristic prediction.
  \item \textbf{Refusal}: the model declines to answer or hedges
  without committing.
\end{itemize}

\noindent The primary metric is the \emph{lure rate}: the proportion of
S1-lure responses on conflict items.  This is more informative than raw
accuracy because it distinguishes heuristic errors from random errors.

% ============================================================
% B. Full Behavioral Results
% ============================================================

\section*{B. Full Behavioral Results}
\label{sec:appendix_behavioral}

\subsection*{B.1 Full Lure Rate Table}

Table~\ref{tab:behavioral_full} extends Table~1 of the main text to
include all categories (including sunk cost where available) and
binomial 95\% confidence intervals.

\begin{table}[ht]
\centering
\caption{S1-lure rates (\%) on conflict items by model and category,
  with binomial exact 95\% confidence intervals in brackets.  Bold
  indicates vulnerability (lure rate $>$10\%).  ``---'' indicates
  category not evaluated for that model.}
\label{tab:behavioral_full}
\vspace{4pt}
\small
\begin{tabular}{lcccc}
\toprule
& \textbf{Llama} & \textbf{R1-Distill} & \multicolumn{2}{c}{\textbf{Qwen 3-8B}} \\
\cmidrule(lr){4-5}
\textbf{Category ($n$ conflict)} & \textbf{8B-Inst.} & \textbf{Llama-8B} & \textbf{no think} & \textbf{think} \\
\midrule
\multicolumn{5}{l}{\emph{Vulnerable categories}} \\
\addlinespace
Base rate ($n$=35) &
  \textbf{84} {\scriptsize [68, 94]} &
  4 {\scriptsize [0, 13]} &
  \textbf{56} {\scriptsize [38, 72]} &
  4 {\scriptsize [0, 13]} \\
\addlinespace
Conjunction ($n$=20) &
  \textbf{55} {\scriptsize [32, 77]} &
  0 {\scriptsize [0, 17]} &
  \textbf{95} {\scriptsize [75, 100]} &
  \textbf{55} {\scriptsize [32, 77]} \\
\addlinespace
Syllogism ($n$=25) &
  \textbf{52} {\scriptsize [31, 72]} &
  0 {\scriptsize [0, 14]} &
  0 {\scriptsize [0, 14]} &
  0 {\scriptsize [0, 14]} \\
\addlinespace
\midrule
\multicolumn{5}{l}{\emph{Immune categories}} \\
\addlinespace
CRT variants ($n$=30) &
  0 {\scriptsize [0, 12]} &
  0 {\scriptsize [0, 12]} &
  0 {\scriptsize [0, 12]} &
  0 {\scriptsize [0, 12]} \\
\addlinespace
Arithmetic ($n$=25) &
  0 {\scriptsize [0, 14]} &
  0 {\scriptsize [0, 14]} &
  0 {\scriptsize [0, 14]} &
  0 {\scriptsize [0, 14]} \\
\addlinespace
Framing ($n$=20) &
  0 {\scriptsize [0, 17]} &
  0 {\scriptsize [0, 17]} &
  0 {\scriptsize [0, 17]} &
  0 {\scriptsize [0, 17]} \\
\addlinespace
Anchoring ($n$=20) &
  0 {\scriptsize [0, 17]} &
  0 {\scriptsize [0, 17]} &
  0 {\scriptsize [0, 17]} &
  0 {\scriptsize [0, 17]} \\
\addlinespace
\midrule
\multicolumn{5}{l}{\emph{Extended categories}} \\
\addlinespace
Sunk cost ($n$=15) &
  --- &
  --- &
  --- &
  --- \\
\addlinespace
\midrule
\textbf{Overall (7-cat)} &
  \textbf{27.3} &
  \textbf{2.4} &
  \textbf{21.0} &
  \textbf{7.0} \\
\bottomrule
\end{tabular}
\end{table}

\subsection*{B.2 Qwen THINK vs.\ NO\_THINK Detailed Comparison}

Table~\ref{tab:qwen_detail} provides a detailed breakdown of the
within-model comparison for Qwen~3-8B.

\begin{table}[ht]
\centering
\caption{Qwen~3-8B THINK vs.\ NO\_THINK: per-category lure rates,
  accuracy, and response classification.  The thinking toggle operates
  on identical weights via the \texttt{/think} and
  \texttt{/no\_think} system prompt flags.}
\label{tab:qwen_detail}
\vspace{4pt}
\small
\begin{tabular}{lcccccc}
\toprule
& \multicolumn{3}{c}{\textbf{NO\_THINK}} & \multicolumn{3}{c}{\textbf{THINK}} \\
\cmidrule(lr){2-4}\cmidrule(lr){5-7}
\textbf{Category} & \textbf{Lure\%} & \textbf{Correct\%} & \textbf{Other\%} & \textbf{Lure\%} & \textbf{Correct\%} & \textbf{Other\%} \\
\midrule
Base rate    & 56 & 44 & 0  & 4  & 96 & 0  \\
Conjunction  & 95 & 5  & 0  & 55 & 45 & 0  \\
Syllogism    & 0  & 48 & 52 & 0  & 100 & 0  \\
CRT variants & 0  & 73 & 27 & 0  & 77 & 23 \\
Arithmetic   & 0  & 100 & 0  & 0  & 96 & 4  \\
Framing      & 0  & 100 & 0  & 0  & 100 & 0  \\
Anchoring    & 0  & 80 & 20 & 0  & 75 & 25 \\
\midrule
\textbf{Overall} & \textbf{21.0} & \textbf{65} & \textbf{14} & \textbf{7.0} & \textbf{85} & \textbf{8} \\
\bottomrule
\end{tabular}
\end{table}

\noindent Key observations from the Qwen within-model comparison:
\begin{itemize}[leftmargin=*,itemsep=2pt]
  \item Base rate neglect shows near-complete correction: 56\% $\to$
  4\% (14$\times$ reduction).
  \item Conjunction fallacy remains majority-susceptible: 95\% $\to$
  55\% (only 1.7$\times$ reduction).  Thinking mode helps but does not
  overcome this deeply encoded heuristic.
  \item Syllogisms are already immune without thinking, unlike
  Llama-3.1-8B where syllogisms are highly vulnerable (52\%).  This
  confirms that vulnerability profiles are model-specific, not
  universal.
  \item All four immune categories remain at 0\% in both modes.
\end{itemize}

% ============================================================
% C. Layer-wise Probe Results (All Layers)
% ============================================================

\section*{C. Layer-Wise Probe Results (All Layers)}
\label{sec:appendix_probes}

\subsection*{C.1 Llama-3.1-8B-Instruct and R1-Distill-Llama-8B (Layers 0--31)}

Table~\ref{tab:probe_all_layers} reports the full layer-wise probe
AUC for both Llama-family models across all 32 layers.  Probes are
$\ell_2$-regularized logistic regression (5-fold stratified CV)
trained on conflict vs.\ no-conflict classification restricted to the
three vulnerable categories (base rate neglect, conjunction fallacy,
syllogisms).

\begin{longtable}{rcccccc}
\caption{Layer-wise probe AUC for Llama-3.1-8B-Instruct and
  R1-Distill-Llama-8B.  ``Gap'' = Llama AUC $-$ R1 AUC.
  Hewitt--Liang control AUC and selectivity (true $-$ control) are
  shown where available.  All values are ROC-AUC from 5-fold
  stratified cross-validation on the three vulnerable categories
  ($n$ = 142 conflict + 142 control items).}
\label{tab:probe_all_layers} \\
\toprule
& \multicolumn{2}{c}{\textbf{Llama 8B-Inst.}} & \multicolumn{2}{c}{\textbf{R1-Distill-Llama-8B}} & & \\
\cmidrule(lr){2-3}\cmidrule(lr){4-5}
\textbf{Layer} & \textbf{AUC} & \textbf{Control} & \textbf{AUC} & \textbf{Control} & \textbf{Gap} & \textbf{Selectivity} \\
\midrule
\endfirsthead
\toprule
\textbf{Layer} & \textbf{AUC} & \textbf{Control} & \textbf{AUC} & \textbf{Control} & \textbf{Gap} & \textbf{Selectivity} \\
\midrule
\endhead
\bottomrule
\endfoot
 0 & 0.848 & 0.605 & 0.782 & 0.534 & 0.066 & 0.299 \\
 1 & 0.899 & 0.586 & 0.831 & 0.568 & 0.068 & 0.317 \\
 2 & 0.926 & 0.603 & 0.870 & 0.584 & 0.056 & 0.290 \\
 3 & 0.932 & 0.588 & 0.901 & 0.565 & 0.031 & 0.337 \\
 4 & 0.927 & 0.610 & 0.919 & 0.555 & 0.008 & 0.320 \\
 5 & 0.939 & 0.626 & 0.928 & 0.599 & 0.011 & 0.308 \\
 6 & 0.933 & 0.601 & 0.926 & 0.570 & 0.007 & 0.334 \\
 7 & 0.960 & 0.610 & 0.927 & 0.560 & 0.033 & 0.331 \\
 8 & 0.952 & 0.603 & 0.920 & 0.567 & 0.032 & 0.302 \\
 9 & 0.948 & 0.549 & 0.928 & 0.577 & 0.020 & 0.388 \\
10 & 0.971 & 0.581 & 0.927 & 0.571 & 0.044 & 0.368 \\
11 & 0.979 & 0.615 & 0.922 & 0.592 & 0.057 & 0.336 \\
12 & 0.988 & 0.578 & 0.927 & 0.572 & 0.061 & 0.370 \\
13 & 0.993 & 0.630 & 0.927 & 0.561 & 0.066 & 0.329 \\
\textbf{14} & \textbf{0.999} & 0.612 & \textbf{0.929} & 0.565 & \textbf{0.070} & 0.347 \\
15 & 0.996 & 0.635 & 0.919 & 0.585 & 0.077 & 0.339 \\
16 & 0.995 & 0.632 & 0.929 & 0.585 & 0.066 & 0.343 \\
17 & 0.989 & 0.596 & 0.923 & 0.567 & 0.066 & 0.372 \\
18 & 0.991 & 0.612 & 0.927 & 0.580 & 0.064 & 0.356 \\
19 & 0.987 & 0.606 & 0.914 & 0.589 & 0.073 & 0.360 \\
20 & 0.978 & 0.626 & 0.914 & 0.589 & 0.064 & 0.333 \\
21 & 0.969 & 0.583 & 0.923 & 0.591 & 0.046 & 0.369 \\
22 & 0.968 & 0.580 & 0.922 & 0.605 & 0.046 & 0.369 \\
23 & 0.963 & 0.600 & 0.917 & 0.585 & 0.046 & 0.348 \\
24 & 0.978 & 0.583 & 0.924 & 0.605 & 0.054 & 0.358 \\
25 & 0.964 & 0.590 & 0.922 & 0.599 & 0.042 & 0.371 \\
26 & 0.965 & 0.580 & 0.918 & 0.594 & 0.047 & 0.358 \\
27 & 0.969 & 0.590 & 0.906 & 0.578 & 0.063 & 0.375 \\
28 & 0.981 & 0.570 & 0.923 & 0.585 & 0.058 & 0.379 \\
29 & 0.965 & 0.572 & 0.924 & 0.590 & 0.041 & 0.362 \\
30 & 0.973 & 0.568 & 0.918 & 0.577 & 0.055 & 0.378 \\
31 & 0.954 & 0.586 & 0.904 & 0.553 & 0.050 & 0.349 \\
\midrule
\textbf{Mean} & \textbf{0.962} & --- & \textbf{0.912} & --- & \textbf{0.050} & --- \\
\textbf{Std} & 0.032 & --- & 0.031 & --- & 0.019 & --- \\
\end{longtable}

\noindent \textbf{Note:} Control AUC columns report Hewitt--Liang probes
trained on randomly permuted labels.  Selectivity = true AUC $-$ control AUC.
The main text
reports aggregate control AUC $\leq 0.64$ at all layers for both
models.  Control AUC ranges from 0.534 to 0.635 for Llama and 0.534
to 0.605 for R1-Distill, yielding selectivity $\geq$29 percentage
points at every layer and $\geq$34 pp at peak layers.

\subsection*{C.2 Summary Statistics by Layer Phase}

\begin{table}[ht]
\centering
\caption{Phase-wise summary of probe AUC.  Phases defined as: early
  (L0--7), mid (L8--15), late (L16--23), final (L24--31).  Gap =
  Llama $-$ R1.  Cohen's $d$ and permutation $p$-value (10{,}000
  permutations, BH-FDR corrected) reported for each phase.}
\label{tab:phase_summary}
\vspace{4pt}
\small
\begin{tabular}{lccccccc}
\toprule
\textbf{Phase} & \textbf{Layers} &
  \textbf{Llama $\bar{\mathrm{AUC}}$} & \textbf{R1 $\bar{\mathrm{AUC}}$} &
  \textbf{Gap} & \textbf{$d$} & \textbf{$p$} & \textbf{Sig.} \\
\midrule
Early & 0--7  & 0.921 & 0.886 & 0.035 & 0.78 & 0.070 & --- \\
Mid   & 8--15 & 0.978 & 0.925 & 0.053 & 3.76 & $<$0.001 & * \\
Late  & 16--23& 0.980 & 0.921 & 0.059 & 6.21 & $<$0.001 & * \\
Final & 24--31& 0.969 & 0.917 & 0.051 & 6.15 & $<$0.001 & * \\
\bottomrule
\end{tabular}
\end{table}

\noindent The deliberation gap widens from early (0.035) to late
(0.059) layers, peaking in the late phase, then narrows slightly in
final layers (0.051).  The early-layer gap does not survive BH-FDR
correction ($p = 0.070$), while mid, late, and final phases are all
significant ($p < 0.001$).

\subsection*{C.3 Qwen~3-8B NO\_THINK and THINK (Layers 0--35)}

Table~\ref{tab:qwen_probe_layers} reports layer-wise probe AUC for
Qwen~3-8B in both thinking modes.  Qwen~3-8B has 36 transformer
layers (layers 0--35).

\begin{longtable}{rccc}
\caption{Layer-wise probe AUC for Qwen~3-8B in NO\_THINK and THINK
  modes.  Bold indicates the peak layer.  Both modes produce identical
  pre-generation representations, reflected in identical probe AUC at
  every layer.}
\label{tab:qwen_probe_layers} \\
\toprule
\textbf{Layer} & \textbf{NO\_THINK AUC} & \textbf{THINK AUC} & \textbf{Difference} \\
\midrule
\endfirsthead
\toprule
\textbf{Layer} & \textbf{NO\_THINK AUC} & \textbf{THINK AUC} & \textbf{Difference} \\
\midrule
\endhead
\bottomrule
\endfoot
 0 & 0.890 & 0.890 & 0.000 \\
 1 & 0.862 & 0.862 & 0.000 \\
 2 & 0.884 & 0.884 & 0.000 \\
 3 & 0.905 & 0.905 & 0.000 \\
 4 & 0.931 & 0.931 & 0.000 \\
 5 & 0.931 & 0.931 & 0.000 \\
 6 & 0.935 & 0.935 & 0.000 \\
 7 & 0.912 & 0.912 & 0.000 \\
 8 & 0.916 & 0.916 & 0.000 \\
 9 & 0.928 & 0.928 & 0.000 \\
10 & 0.930 & 0.930 & 0.000 \\
11 & 0.924 & 0.924 & 0.000 \\
12 & 0.924 & 0.924 & 0.000 \\
13 & 0.936 & 0.936 & 0.000 \\
14 & 0.921 & 0.921 & 0.000 \\
15 & 0.924 & 0.924 & 0.000 \\
16 & 0.937 & 0.937 & 0.000 \\
17 & 0.938 & 0.938 & 0.000 \\
18 & 0.948 & 0.948 & 0.000 \\
19 & 0.957 & 0.957 & 0.000 \\
20 & 0.965 & 0.965 & 0.000 \\
21 & 0.965 & 0.965 & 0.000 \\
22 & 0.969 & 0.969 & 0.000 \\
23 & 0.965 & 0.965 & 0.000 \\
24 & 0.969 & 0.969 & 0.000 \\
25 & 0.967 & 0.967 & 0.000 \\
26 & 0.951 & 0.951 & 0.000 \\
27 & 0.960 & 0.960 & 0.000 \\
28 & 0.957 & 0.957 & 0.000 \\
29 & 0.963 & 0.963 & 0.000 \\
30 & 0.963 & 0.963 & 0.000 \\
31 & 0.964 & 0.964 & 0.000 \\
32 & 0.945 & 0.945 & 0.000 \\
33 & 0.966 & 0.966 & 0.000 \\
\textbf{34} & \textbf{0.971} & \textbf{0.971} & \textbf{0.000} \\
35 & 0.959 & 0.959 & 0.000 \\
\end{longtable}

\noindent \textbf{Note:} The think/no-think difference is exactly
0.000 at every layer because both modes share identical weights and
the \pzero{} position is deterministic given the prompt.  The thinking
toggle only affects generation behavior \emph{after} \pzero{}.  Per-layer
Qwen peak AUC is 0.971 at L34 on the vulnerable-category subset and
0.997 at L24 on all categories.  The peak layer for vulnerable categories
shifts to L24 (tied at 0.969 with L22) on the vulnerable subset.
Hewitt--Liang control probes for Qwen are pending.

% ============================================================
% D. Cross-Prediction Details
% ============================================================

\section*{D. Cross-Prediction Details}
\label{sec:appendix_cross}

Cross-prediction tests train a probe on the three \emph{vulnerable}
categories (base rate neglect, conjunction fallacy, syllogisms) and
evaluate it on the four \emph{immune} categories (CRT, arithmetic,
framing, anchoring).  If the probe captures a universal
processing-mode signal, transfer AUC should remain high; if it
captures category-specific lure susceptibility, transfer AUC should
drop to chance ($\approx$0.5) or below.

\subsection*{D.1 Full Cross-Prediction Results}

Table~\ref{tab:cross_pred_full} reports cross-prediction AUC at all
10 tested layers for both Llama-family models.

\begin{table}[ht]
\centering
\caption{Cross-prediction AUC at tested layers.  ``Within'' = probe
  AUC on vulnerable categories (same-category, different-fold).
  ``Transfer'' = probe AUC on immune categories.  ``Drop'' =
  Within $-$ Transfer.}
\label{tab:cross_pred_full}
\vspace{4pt}
\small
\begin{tabular}{rcccccc}
\toprule
& \multicolumn{3}{c}{\textbf{Llama 8B-Inst.}} & \multicolumn{3}{c}{\textbf{R1-Distill-Llama-8B}} \\
\cmidrule(lr){2-4}\cmidrule(lr){5-7}
\textbf{Layer} & \textbf{Within} & \textbf{Transfer} & \textbf{Drop} & \textbf{Within} & \textbf{Transfer} & \textbf{Drop} \\
\midrule
 0 & 0.848 & 0.702 & 0.146 & 0.782 & 0.734 & 0.048 \\
 4 & 0.927 & 0.439 & 0.488 & 0.919 & 0.878 & 0.041 \\
 8 & 0.952 & 0.384 & 0.568 & 0.920 & 0.878 & 0.042 \\
12 & 0.988 & 0.449 & 0.539 & 0.927 & 0.728 & 0.199 \\
\textbf{14} & \textbf{0.999} & \textbf{0.378} & \textbf{0.621} & \textbf{0.929} & \textbf{0.685} & \textbf{0.244} \\
16 & 0.995 & 0.569 & 0.426 & 0.929 & 0.696 & 0.233 \\
20 & 0.978 & 0.386 & 0.592 & 0.914 & 0.615 & 0.299 \\
24 & 0.978 & 0.326 & 0.652 & 0.924 & 0.524 & 0.400 \\
28 & 0.981 & 0.365 & 0.616 & 0.923 & 0.415 & 0.508 \\
31 & 0.954 & 0.438 & 0.516 & 0.904 & 0.385 & 0.519 \\
\midrule
\textbf{Mean} & 0.960 & 0.444 & 0.516 & 0.907 & 0.654 & 0.253 \\
\bottomrule
\end{tabular}
\end{table}

\subsection*{D.2 Interpretation}

\paragraph{Llama-3.1-8B-Instruct.}
Transfer AUC drops dramatically from within-category performance at
all mid-to-late layers.  At the peak layer (L16), transfer AUC is
0.378---\emph{below} chance (0.5), indicating that the probe's
learned direction is \emph{anticorrelated} with the immune-category
conflict/control distinction.  This anticorrelation is strong
evidence that the probe captures processing-mode information specific
to the vulnerable categories, not generic surface features of lure
text.  The mean transfer AUC across all tested layers is 0.444
(bootstrap 95\% CI: [0.386, 0.515]).

\paragraph{R1-Distill-Llama-8B.}
The pattern is qualitatively different.  At early layers (L0--L8),
transfer AUC remains high ($\geq$0.734), suggesting that R1-Distill's
early representations encode a more generic conflict/control signal
that partially transfers across categories.  Transfer AUC declines
monotonically in later layers, reaching 0.385 at L31---comparable to
Llama's specificity.  The mean transfer AUC is 0.654 (bootstrap 95\%
CI: [0.551, 0.752]), significantly above Llama's ($p < 0.05$, paired
bootstrap test).

\paragraph{Layer-wise specificity trajectory.}
For Llama, specificity is established by L4 (drop = 0.488) and
remains high thereafter.  For R1-Distill, specificity develops
gradually, with the drop exceeding 0.3 only from L20 onward.  This
suggests that reasoning distillation delays category-specific encoding
to later layers while maintaining a more generic early representation.

% ============================================================
% E. Transfer Matrix
% ============================================================

\section*{E. Transfer Matrix}
\label{sec:appendix_transfer}

The transfer matrix tests whether probes trained on individual
categories transfer to other categories, revealing shared vs.\
distinct representational structure across bias types.  We train a
logistic regression probe on each category independently and test on
all other categories at layer 14 of Llama-3.1-8B-Instruct (the peak
probing layer).

\subsection*{E.1 Full Transfer Matrix (Llama L14)}

Table~\ref{tab:transfer_matrix} reports the full category-to-category
transfer AUC matrix.

\begin{table}[ht]
\centering
\caption{Category transfer matrix (Llama-3.1-8B-Instruct, L14).
  Rows = training category; columns = test category.  Diagonal
  entries are within-category 5-fold CV AUC.  Off-diagonal entries
  are transfer AUC (probe trained on row category, tested on column
  category).  ``Vuln.'' = vulnerable categories; ``Imm.'' = immune
  categories.}
\label{tab:transfer_matrix}
\vspace{4pt}
\small
\begin{tabular}{lccc|cccc}
\toprule
& \multicolumn{3}{c}{\textbf{Vuln.\ (test)}} & \multicolumn{4}{c}{\textbf{Imm.\ (test)}} \\
\cmidrule(lr){2-4}\cmidrule(lr){5-8}
\textbf{Train} & \textbf{Base} & \textbf{Conj.} & \textbf{Syll.} & \textbf{CRT} & \textbf{Arith.} & \textbf{Fram.} & \textbf{Anch.} \\
\midrule
Base rate    & 1.000 & 0.993 & 0.594 & 0.491 & 0.954 & 0.370 & 0.440 \\
Conjunction  & 0.998 & 1.000 & 0.627 & 0.334 & 0.000 & 0.000 & 0.787 \\
Syllogism    & 0.950 & 0.990 & 0.936 & 0.457 & 0.000 & 0.875 & 0.000 \\
\midrule
CRT          & 0.062 & 0.223 & 0.451 & 1.000 & 0.000 & 0.000 & 0.995 \\
Arithmetic   & 0.634 & 0.062 & 0.442 & 0.453 & 1.000 & 0.815 & 0.580 \\
Framing      & 0.539 & 0.057 & 0.571 & 0.447 & 0.594 & 1.000 & 0.108 \\
Anchoring    & 0.378 & 0.648 & 0.190 & 0.540 & 0.906 & 0.000 & 1.000 \\
\bottomrule
\end{tabular}
\end{table}

\subsection*{E.2 Interpretation}

The partial transfer matrix reveals a structured pattern.  Key
findings from the available data:

\begin{itemize}[leftmargin=*,itemsep=4pt]
  \item \textbf{Base rate $\leftrightarrow$ conjunction transfer is
  near-perfect} (AUC = 0.993 in both directions).  These two
  representativeness-heuristic categories share a processing-mode
  direction that is nearly identical in the residual stream at L16.

  \item \textbf{CRT probes do not transfer to vulnerable categories}
  (AUC = 0.062, as reported in the main text).  This confirms that
  CRT items, where all models are immune, occupy a distinct region of
  representational space.

  \item The high within-vulnerable transfer (AUC $\geq$ 0.950 among
  base rate, conjunction, and syllogism) combined with the low
  CRT-to-vulnerable transfer (0.062) provides converging evidence
  that the probe captures a shared processing-mode signal specific to
  the vulnerable categories, not a generic surface-feature signal.

  \item \textbf{Syllogism transfers well to base rate and conjunction}
  (AUC = 0.950 and 0.990 respectively), completing the
  within-vulnerable transfer triangle.

  \item \textbf{Immune categories show heterogeneous cross-transfer.}
  CRT$\to$base rate is 0.062 (near anti-correlated), while
  anchoring$\to$arithmetic reaches 0.906.  The immune categories do
  not share a single processing-mode direction.
\end{itemize}

% ============================================================
% F. Representational Geometry
% ============================================================

\section*{F. Representational Geometry}
\label{sec:appendix_geometry}

We assess the geometric structure of conflict vs.\ no-conflict
representations using two complementary metrics: cosine silhouette
scores (cluster separability) and Centered Kernel Alignment (CKA;
representational similarity between models).

\subsection*{F.1 Silhouette Scores}

Silhouette scores measure how well conflict and no-conflict items
cluster in residual stream space.  We compute cosine silhouette after
PCA reduction to 100 dimensions (to address the $d \gg N$ pitfall
per Cover's theorem).

\begin{table}[ht]
\centering
\caption{Cosine silhouette scores by layer phase.  Positive values
  indicate meaningful clustering of conflict vs.\ no-conflict items.
  Values near zero indicate overlapping clusters.  All values computed
  after PCA to 100 dimensions.}
\label{tab:silhouette}
\vspace{4pt}
\small
\begin{tabular}{lccccc}
\toprule
\textbf{Model} & \textbf{Peak Sil.} & \textbf{Peak Layer} &
  \textbf{Mean Sil.\ (L0--31)} & \textbf{Min Sil.} & \textbf{Min Layer} \\
\midrule
Llama 8B-Inst. & 0.079 & 13 & 0.037 & 0.009 & 0 \\
R1-Distill     & 0.059 & 0  & 0.044 & 0.032 & 1 \\
\bottomrule
\end{tabular}
\end{table}

\noindent Both models show positive but low silhouette scores,
indicating that conflict and no-conflict items form overlapping
clusters with only modest geometric separation.  Llama's higher
silhouette (0.079 vs.\ 0.059) is consistent with its higher probe
AUC---the linear boundary is more distinct when clusters are
marginally better separated.  However, the low absolute values
(far below the 0.25 threshold for ``reasonable structure'') suggest
that the conflict/control distinction is a subtle direction in high-dimensional
space rather than a gross geometric partition.

\paragraph{Layer-wise silhouette profiles.}
Full layer-wise silhouette scores for both models are reported in
Table~\ref{tab:silhouette_layers}.

\begin{longtable}{rcc}
\caption{Layer-wise cosine silhouette scores (PCA to 100 dimensions).}
\label{tab:silhouette_layers} \\
\toprule
\textbf{Layer} & \textbf{Llama 8B-Inst.} & \textbf{R1-Distill-Llama-8B} \\
\midrule
\endfirsthead
\toprule
\textbf{Layer} & \textbf{Llama 8B-Inst.} & \textbf{R1-Distill-Llama-8B} \\
\midrule
\endhead
\bottomrule
\endfoot
 0 & 0.009 & 0.059 \\
 1 & 0.012 & 0.032 \\
 2 & 0.022 & 0.039 \\
 3 & 0.030 & 0.039 \\
 4 & 0.026 & 0.046 \\
 5 & 0.029 & 0.043 \\
 6 & 0.035 & 0.040 \\
 7 & 0.043 & 0.042 \\
 8 & 0.049 & 0.040 \\
 9 & 0.042 & 0.041 \\
10 & 0.051 & 0.043 \\
11 & 0.054 & 0.044 \\
12 & 0.058 & 0.047 \\
13 & 0.079 & 0.051 \\
14 & 0.077 & 0.053 \\
15 & 0.074 & 0.052 \\
16 & 0.061 & 0.048 \\
17 & 0.057 & 0.046 \\
18 & 0.055 & 0.046 \\
19 & 0.050 & 0.045 \\
20 & 0.041 & 0.045 \\
21 & 0.038 & 0.045 \\
22 & 0.031 & 0.043 \\
23 & 0.029 & 0.043 \\
24 & 0.023 & 0.041 \\
25 & 0.020 & 0.041 \\
26 & 0.017 & 0.042 \\
27 & 0.014 & 0.042 \\
28 & 0.015 & 0.042 \\
29 & 0.015 & 0.042 \\
30 & 0.012 & 0.041 \\
31 & 0.009 & 0.041 \\
\end{longtable}

\subsection*{F.2 Centered Kernel Alignment (CKA)}

CKA \citep{kornblith2019similarity} measures the similarity of
representational geometry between the two Llama-family models at each
layer.  A CKA value of 1.0 indicates identical representational
structure (up to isotropic scaling); 0.0 indicates no linear
relationship.

\begin{table}[ht]
\centering
\caption{CKA (linear kernel) between Llama-3.1-8B-Instruct and
  R1-Distill-Llama-8B residual stream activations at matched layers.}
\label{tab:cka}
\vspace{4pt}
\small
\begin{tabular}{rcrcrcrcrcrcrcrcrcrcrcrcrcrcrcrcrcrr}
\toprule
\multicolumn{2}{c}{\textbf{Summary}} \\
\midrule
Min CKA & 0.379 \\
Max CKA & 0.985 \\
Mean CKA & 0.929 \\
\bottomrule
\end{tabular}
\end{table}

\paragraph{Layer-wise CKA.}
Full layer-wise CKA values between the two models:

\begin{longtable}{rc}
\caption{Layer-wise CKA (linear kernel) between Llama-3.1-8B-Instruct
  and R1-Distill-Llama-8B.}
\label{tab:cka_layers} \\
\toprule
\textbf{Layer} & \textbf{CKA} \\
\midrule
\endfirsthead
\toprule
\textbf{Layer} & \textbf{CKA} \\
\midrule
\endhead
\bottomrule
\endfoot
 0 & 0.379 \\
 1 & 0.664 \\
 2 & 0.763 \\
 3 & 0.876 \\
 4 & 0.888 \\
 5 & 0.966 \\
 6 & 0.971 \\
 7 & 0.975 \\
 8 & 0.985 \\
 9 & 0.982 \\
10 & 0.962 \\
11 & 0.956 \\
12 & 0.966 \\
13 & 0.971 \\
14 & 0.976 \\
15 & 0.957 \\
16 & 0.970 \\
17 & 0.976 \\
18 & 0.978 \\
19 & 0.975 \\
20 & 0.972 \\
21 & 0.963 \\
22 & 0.971 \\
23 & 0.962 \\
24 & 0.971 \\
25 & 0.969 \\
26 & 0.961 \\
27 & 0.964 \\
28 & 0.964 \\
29 & 0.965 \\
30 & 0.960 \\
31 & 0.955 \\
\end{longtable}

\noindent The CKA range (0.379--0.985) indicates that
reasoning distillation preserves substantial representational structure
at some layers while substantially reorganizing others.  We expect
that layers where CKA is lowest (furthest from 1.0) correspond to
where the deliberation gap in probe AUC is largest, as those are the
layers most affected by reasoning training.  Full layer-wise CKA
values will test this prediction.

% ============================================================
% G. Lure Susceptibility Distribution
% ============================================================

\section*{G. Lure Susceptibility Distribution}
\label{sec:appendix_lure}

Lure susceptibility scores are continuous \pzero{} logit scores for
the S1-lure answer on each conflict item.  Positive scores indicate
that the model's pre-generation representation favors the heuristic
(lure) answer; negative scores indicate it favors the correct
(normative) answer.

\subsection*{G.1 Distribution Summary}

\begin{table}[ht]
\centering
\caption{Lure susceptibility score distributions.  Scores are
  \pzero{} logit differences (lure $-$ correct) on conflict items
  ($n$ = 165 per model, 7 primary categories).}
\label{tab:lure_dist}
\vspace{4pt}
\small
\begin{tabular}{lcccccc}
\toprule
\textbf{Model} & \textbf{Mean} & \textbf{SD} & \textbf{Median} &
  \textbf{Min} & \textbf{Max} & \textbf{$n$ positive} \\
\midrule
Llama 8B-Inst.   & $+$0.422 & 2.995 & 0.000 & $-$6.938 & $+$8.719 & 56 \\
R1-Distill       & $-$0.326 & 2.063 & $-$0.320 & $-$7.301 & $+$5.219 & 73 \\
\midrule
\multicolumn{7}{l}{Welch's $t$(250.2) = 2.459, $p$ = 0.007 (one-sided), Cohen's $d$ = 0.292} \\
\multicolumn{7}{l}{Permutation test: $p$ = 0.008 (10{,}000 permutations)} \\
\bottomrule
\end{tabular}
\end{table}

\noindent The sign flip between models is the key finding: Llama's
pre-generation representation actively favors the lure ($+$0.422),
while R1-Distill's favors the correct answer ($-$0.326).  The effect
size is small (Cohen's $d$ = 0.292), reflecting the high variance
within each model, but is statistically significant by both
parametric and non-parametric tests.

\subsection*{G.2 Per-Category Breakdown}

\begin{table}[ht]
\centering
\caption{Lure susceptibility scores by category and model.
  Positive = model favors lure; negative = model favors correct
  answer.}
\label{tab:lure_by_category}
\vspace{4pt}
\small
\begin{tabular}{lcccc}
\toprule
& \multicolumn{2}{c}{\textbf{Llama 8B-Inst.}} & \multicolumn{2}{c}{\textbf{R1-Distill}} \\
\cmidrule(lr){2-3}\cmidrule(lr){4-5}
\textbf{Category} & \textbf{Mean} & \textbf{SD} & \textbf{Mean} & \textbf{SD} \\
\midrule
Base rate ($n$=25)    & $+$1.421 & 3.985 & $-$2.184 & 3.070 \\
Conjunction ($n$=20)  & $+$3.082 & 1.601 & $-$0.865 & 0.489 \\
Syllogism ($n$=25)    & $-$1.011 & 5.652 & $+$0.068 & 1.992 \\
CRT ($n$=30)          & $+$0.000 & 0.000 & $+$0.394 & 1.527 \\
Arithmetic ($n$=25)   & $+$0.000 & 0.000 & $-$0.359 & 1.878 \\
Framing ($n$=20)      & $-$0.111 & 1.409 & $-$0.127 & 1.053 \\
Anchoring ($n$=20)    & $+$0.000 & 0.000 & $+$0.808 & 1.778 \\
\midrule
\textbf{Overall} ($n$=165) & $+$0.422 & 2.995 & $-$0.326 & 2.063 \\
\bottomrule
\end{tabular}
\end{table}

\noindent We expect the per-category breakdown to show that the
sign flip is driven primarily by the vulnerable categories (base
rate, conjunction, syllogism), where Llama shows high lure rates
and large positive susceptibility scores.  For the immune categories
(CRT, arithmetic, framing, anchoring), both models should show
negative susceptibility scores (favoring the correct answer).  The
per-category data will test whether the overall effect is driven by
a few categories or represents a broad shift.

\subsection*{G.3 Histogram Description}

The lure susceptibility distributions for both models are approximately
normal with heavy tails.  Llama's distribution is centered at $+$0.422
with standard deviation 2.995, indicating that the typical conflict
item produces a moderate pro-lure bias in the pre-generation
representation.  R1-Distill's distribution is centered at $-$0.326
with standard deviation 2.063, shifted toward anti-lure bias with
notably lower variance.  The variance reduction in R1-Distill
($\sigma^2 = 4.25$ vs.\ $8.97$ for Llama, $F$-test $p$
$< 0.001$, $F = 2.11$) suggests that reasoning distillation not only shifts
the mean but also reduces the variability of lure susceptibility
across items.

% ============================================================
% H. OLMo-3 Replication (Probes and Behavioral)
% ============================================================

\section*{H. OLMo-3 Replication}
\label{sec:appendix_olmo}

To test the generality of our findings beyond the Llama architecture
family, we replicate the probing and behavioral analyses on OLMo-3-7B
\citep{groeneveld2024olmo}, using its Instruct and Think variants as a
second controlled reasoning pair (same base weights, reasoning-tuned
vs.\ standard instruction-tuned).

\subsection*{H.1 OLMo-3 Behavioral Results}

Both OLMo-3 variants show \textbf{0\% S1-lure rates across all
categories}, including the three categories where Llama-3.1-8B is
highly vulnerable.  This makes OLMo-3 an important control: it
demonstrates that S1-lure vulnerability is model-specific and not an
inherent property of the benchmark items.

\begin{table}[ht]
\centering
\caption{OLMo-3 S1-lure rates (\%) on conflict items.
  Binomial exact 95\% CIs in brackets.}
\label{tab:olmo_behavioral}
\vspace{4pt}
\small
\begin{tabular}{lcc}
\toprule
\textbf{Category ($n$ conflict)} & \textbf{OLMo-3 Instruct} & \textbf{OLMo-3 Think} \\
\midrule
Base rate ($n$=35)    & 0 {\scriptsize [0, 8]}  & 0 {\scriptsize [0, 8]}  \\
Conjunction ($n$=20)  & 0 {\scriptsize [0, 14]} & 0 {\scriptsize [0, 14]} \\
Syllogism ($n$=25)    & 0 {\scriptsize [0, 11]} & 0 {\scriptsize [0, 11]} \\
CRT ($n$=30)          & 0 {\scriptsize [0, 10]} & 0 {\scriptsize [0, 10]} \\
Arithmetic ($n$=25)   & 0 {\scriptsize [0, 11]} & 0 {\scriptsize [0, 11]} \\
Framing ($n$=20)      & 0 {\scriptsize [0, 14]} & 0 {\scriptsize [0, 14]} \\
Anchoring ($n$=20)    & 0 {\scriptsize [0, 14]} & 0 {\scriptsize [0, 14]} \\
Sunk cost ($n$=15)    & 0 {\scriptsize [0, 18]} & 0 {\scriptsize [0, 18]} \\
\midrule
\textbf{Overall (7-cat)} & \textbf{0.0} & \textbf{0.0} \\
\bottomrule
\end{tabular}
\end{table}

\subsection*{H.2 OLMo-3 Probe Results}

Despite both models achieving 0\% lure rates, probes still
discriminate conflict from control items at high AUC, confirming that
conflict-detection representations exist even in behaviourally immune
models.

\begin{table}[ht]
\centering
\caption{OLMo-3 probe summary (vulnerable categories, \pzero{} position,
  $n$ = 160).  Both models have 32 layers.}
\label{tab:olmo_probe_summary}
\vspace{4pt}
\small
\begin{tabular}{lcccc}
\toprule
\textbf{Model} & \textbf{Peak Layer} & \textbf{Peak AUC} & \textbf{Mean AUC} & \textbf{Min AUC} \\
\midrule
OLMo-3-Instruct & 24 & 0.996 [0.988, 1.000] & 0.969 & 0.861 \\
OLMo-3-Think    & 22 & 0.962 [0.934, 0.982] & 0.946 & 0.841 \\
\midrule
\multicolumn{5}{l}{Instruct higher at 30/32 layers; max gap 0.068 at L12; mean gap 0.023} \\
\bottomrule
\end{tabular}
\end{table}

\noindent Key findings:
\begin{itemize}[leftmargin=*,itemsep=2pt]
  \item \textbf{OLMo-3-Instruct peaks at L24} (AUC = 0.996 [0.988, 1.000]),
  similar to Llama's peak at L16 (in the 60--75\% depth range of the network).
  \item \textbf{OLMo-3-Think peaks at L22} (AUC = 0.962 [0.934, 0.982]),
  with bootstrap CIs non-overlapping with Instruct, confirming the gap
  is statistically robust (0.034 AUC).
  \item The Instruct model has higher probe AUC at 30 of 32 layers,
  mirroring the Llama $>$ R1 pattern despite both OLMo models being
  behaviourally immune.
  \item All P2 (control position) probes return AUC = 0.500 exactly,
  confirming that the signal is position-specific.
\end{itemize}

\paragraph{Layer-wise OLMo-3 probe AUC.}

\begin{longtable}{rcccc}
\caption{Layer-wise probe AUC for OLMo-3-Instruct and OLMo-3-Think
  (vulnerable categories, \pzero{} position, 5-fold CV,
  mean $\pm$ std over folds).}
\label{tab:olmo_probe_layers} \\
\toprule
\textbf{Layer} & \textbf{Instruct AUC} & \textbf{$\pm$ std} & \textbf{Think AUC} & \textbf{$\pm$ std} \\
\midrule
\endfirsthead
\toprule
\textbf{Layer} & \textbf{Instruct AUC} & \textbf{$\pm$ std} & \textbf{Think AUC} & \textbf{$\pm$ std} \\
\midrule
\endhead
\bottomrule
\endfoot
 0 & 0.861 & 0.005 & 0.846 & 0.009 \\
 1 & 0.868 & 0.007 & 0.841 & 0.015 \\
 2 & 0.895 & 0.017 & 0.882 & 0.008 \\
 3 & 0.903 & 0.014 & 0.886 & 0.004 \\
 4 & 0.907 & 0.004 & 0.892 & 0.005 \\
 5 & 0.913 & 0.009 & 0.897 & 0.003 \\
 6 & 0.948 & 0.002 & 0.899 & 0.003 \\
 7 & 0.940 & 0.001 & 0.903 & 0.001 \\
 8 & 0.956 & 0.004 & 0.902 & 0.004 \\
 9 & 0.968 & 0.000 & 0.906 & 0.001 \\
10 & 0.970 & 0.003 & 0.905 & 0.002 \\
11 & 0.969 & 0.005 & 0.906 & 0.005 \\
12 & 0.991 & 0.001 & 0.922 & 0.002 \\
13 & 0.992 & 0.002 & 0.937 & 0.004 \\
14 & 0.995 & 0.001 & 0.947 & 0.006 \\
15 & 0.997 & 0.001 & 0.976 & 0.004 \\
16 & 0.998 & 0.002 & 0.989 & 0.003 \\
17 & 0.997 & 0.001 & 0.988 & 0.004 \\
18 & 0.997 & 0.001 & 0.986 & 0.004 \\
19 & 0.998 & 0.002 & 0.987 & 0.003 \\
20 & 0.998 & 0.001 & 0.986 & 0.003 \\
21 & 0.998 & 0.001 & 0.985 & 0.002 \\
\textbf{22} & 0.998 & 0.001 & \textbf{0.991} & 0.002 \\
23 & 0.998 & 0.001 & 0.992 & 0.002 \\
\textbf{24} & \textbf{0.998} & 0.001 & 0.992 & 0.002 \\
25 & 0.997 & 0.001 & 0.991 & 0.002 \\
26 & 0.997 & 0.002 & 0.992 & 0.002 \\
27 & 0.996 & 0.004 & 0.992 & 0.003 \\
28 & 0.994 & 0.006 & 0.993 & 0.002 \\
29 & 0.994 & 0.005 & 0.992 & 0.002 \\
30 & 0.990 & 0.007 & 0.991 & 0.002 \\
31 & 0.986 & 0.006 & 0.991 & 0.001 \\
\end{longtable}

% ============================================================
% I. SAE Feature Analysis
% ============================================================

\section*{I. SAE Feature Analysis}
\label{sec:appendix_sae}

We use the Goodfire sparse autoencoder (SAE) trained on
Llama-3.1-8B-Instruct layer~19 to decompose the residual stream
into interpretable features and test which features are associated
with the conflict/control distinction.

\subsection*{I.1 Reconstruction Fidelity}

\begin{table}[ht]
\centering
\caption{SAE reconstruction fidelity on Llama-3.1-8B-Instruct
  layer~19 activations ($n$ = 470 items).}
\label{tab:sae_fidelity}
\vspace{4pt}
\small
\begin{tabular}{lc}
\toprule
\textbf{Metric} & \textbf{Value} \\
\midrule
MSE                  & 0.011 \\
Explained variance   & 74.0\% \\
Mean L0 (active features) & 47.9 \\
SAE width            & 65{,}536 \\
Is poor fit?         & No \\
\bottomrule
\end{tabular}
\end{table}

\noindent The explained variance of 74\% is moderate, indicating that
the SAE captures the majority of variance but not all.  Downstream
feature-level findings should be interpreted as a lower bound on
the true signal, since reconstruction loss may attenuate feature
activations.

\subsection*{I.2 Significant Features}

After BH-FDR correction ($q < 0.05$) across 65{,}536 features, we
identify \textbf{41 features} with significantly different activation
between conflict and control items.

\subsection*{I.3 Ma et al.\ Falsification}

All 41 significant features were submitted to the Ma et al.\ (2026)
falsification protocol: for each feature, the top-3 activating tokens
were injected into 100 random non-cognitive-bias texts.  A feature
is declared spurious if it still activates above threshold in the
random context.

\begin{table}[ht]
\centering
\caption{SAE falsification results (Llama-3.1-8B-Instruct, L19).}
\label{tab:sae_falsification}
\vspace{4pt}
\small
\begin{tabular}{lc}
\toprule
\textbf{Metric} & \textbf{Value} \\
\midrule
Features tested (post-FDR)      & 41 \\
Falsification candidates        & 41 \\
Spurious (failed falsification) & 0 \\
Surviving features              & 41 \\
Spurious rate                   & 0.0\% \\
\bottomrule
\end{tabular}
\end{table}

\noindent All 41 features survive falsification.  The 0\% spurious
rate indicates that none of the significant features are simple
token-level artifacts --- they activate specifically in the context
of cognitive bias items, not when the same tokens appear in random
text.

% ============================================================
% J. Attention Entropy Analysis
% ============================================================

\section*{J. Attention Entropy Analysis}
\label{sec:appendix_attention}

We compare attention entropy patterns between Llama-3.1-8B-Instruct
and R1-Distill-Llama-8B at the \pzero{} position.  Attention entropy
is computed as the Shannon entropy of each head's attention
distribution, normalized by $\log_2(t)$ where $t$ is the sequence
length.  Head-level significance is assessed via rank-biserial
correlation between conflict and control items, with BH-FDR
correction at $q < 0.05$.

\subsection*{J.1 Head-Level Results}

\begin{table}[ht]
\centering
\caption{Summary of head-level attention entropy tests.}
\label{tab:attention_heads}
\vspace{4pt}
\small
\begin{tabular}{lcc}
\toprule
\textbf{Metric} & \textbf{Llama 8B-Inst.} & \textbf{R1-Distill} \\
\midrule
Total heads tested                 & 1{,}024 & 1{,}024 \\
Significant (BH $q < 0.05$)       & 473 (46.2\%) & 515 (50.3\%) \\
\quad Of which conflict-higher     & 82 (17.3\%) & 115 (22.3\%) \\
\quad Of which control-higher      & 391 (82.7\%) & 400 (77.7\%) \\
\midrule
S2-specialized heads               & 30 (2.9\%)  & 57 (5.6\%) \\
\quad ($q < 0.05$ and $|r_{rb}| \geq 0.3$) & & \\
\quad Of which conflict-higher     & 29 (96.7\%) & 47 (82.5\%) \\
\quad KV-group level               & 13 (5.1\%)  & 24 (9.4\%) \\
\bottomrule
\end{tabular}
\end{table}

\subsection*{J.2 S2-Specialized Heads}

S2-specialized heads are defined as those with both statistical
significance ($q < 0.05$) and a meaningful effect size ($|r_{rb}|
\geq 0.3$) on the normalized entropy metric.  R1-Distill has
nearly twice the proportion of S2-specialized heads compared to
Llama (5.6\% vs.\ 2.9\%), and this gap widens further at the
KV-group level (9.4\% vs.\ 5.1\%), accounting for the
non-independence of query heads within GQA groups.

\subsection*{J.3 Layer-Wise Entropy Gap}

The entropy gap (conflict $-$ control mean normalized entropy) shows
distinct layer profiles for each model:

\begin{itemize}[leftmargin=*,itemsep=2pt]
  \item \textbf{Llama-3.1-8B:} Peak entropy gap at layer~7
  ($\Delta = -0.0156$), concentrated in early-to-mid layers
  (L4--L14).
  \item \textbf{R1-Distill:} Peak entropy gap at layer~11
  ($\Delta = -0.0167$), with a broader distribution extending
  through layer~15 and a secondary focus in late layers (L27--L31).
\end{itemize}

\noindent The negative sign indicates that conflict items elicit
\emph{lower} entropy (more focused attention) than control items,
consistent with the hypothesis that the heuristic--normative conflict
triggers a more targeted attention pattern.

\vspace{1em}
\hrule
\vspace{0.5em}
\noindent\small\textbf{Data availability.}  All result JSON files,
analysis scripts, and benchmark items are available at the project
repository.

% ============================================================
% References
% ============================================================

\bibliographystyle{plainnat}
\bibliography{references}

\end{document}
